% !TeX spellcheck = it_IT

\documentclass[aspectratio=169, handout]{beamer}

\usepackage[utf8]{inputenc}
\usepackage[italian]{babel}
\usepackage{multicol}

\usetheme[titleformat=smallcaps, numbering=fraction, background=light, progressbar=frametitle]{metropolis}

\title{Esercitazioni di Informatica B}
\subtitle{Struct}

\author{Stefano Cereda\\
	stefano.cereda@polimi.it
}
\date{27/10/2018}
\institute[PoliMi]{Politecnico Milano}
\logo{\includegraphics[width=10mm]{../logopolimi}}

\setbeamercovered{invisible}

\makeindex

\begin{document}
	\begin{frame}
	\maketitle
\end{frame}

%\begin{frame}<handout:0>{Nuovi esercizi su argomenti vecchi}
%Aggiunti due esercizi sulle matrici e alcuni esercizi sul binario.
%\end{frame}

\begin{frame}{Struct e matrici}
    Si definisca una struttura dati \emph{eta\_nome} per contenere un numero ed una stringa.
    Si crei poi una matrice $N x N$ di \emph{eta\_nome}.

    Si scriva un programma che chieda all'utente di riempire questa matrice, indicando nome ed età delle persone sedute
    in un'aula in base alla loro posizione.

    Infine, si stampino i nomi delle persone maggiorenni sedute sulla diagonale dell'aula.
\end{frame}

\begin{frame}[allowframebreaks]{Scontrini}
Un supermercato ha memorizzato il proprio archivio di scontrini nell'array \emph{archivio}, in cui ogni elemento è una struttura dati di tipo \emph{scontrino} con le  seguenti informazioni:
\begin{itemize}
	\item IDcliente: stringa di 10 caratteri che identifica univocamente il cliente
	\item totale: totale della spesa in Euro
	\item punti: punti premio associati alla spesa
\end{itemize}
Ad ogni spesa, vengono assegnati ad ogni cliente un quantitativo di punti premio pari alla somma dei punti raccolti più un ulteriore punto premio per ogni 10 euro spesi.

\begin{enumerate}
	\item  Definire il tipo $scontrino$ e l'array $archivio$
	\item Scrivere il frammento di codice C che, dato il codice identificativo di un cliente, calcoli la somma dei suoi punti premio.
	\item Scrivere il frammento di codice C che ordina l'array \emph{archivio} secondo il totale della spesa in ordine decrescente.
\end{enumerate}
Nota: Si faccia attenzione al fatto che un cliente può comparire in più di uno scontrino.
\end{frame}


\begin{frame}{Fantacalcio da TdE 17/09/2012}
Definire un tipo di dato \emph{fantacalcio} che contiene una lista di $N = 20$ \emph{fantasquadre}. Ogni fantasquadra è identificata dal nome, il colore della casacca, il punteggio corrente e l’allenatore. L’\emph{allenatore} è definito dal nome, cognome e numero di coppe vinte (un intero $\geq 0$). Infine si dichiari una variabile \emph{campionato} di
tipo fantacalcio.

Scrivere un frammento di codice in linguaggio C che stampi a video il numero totale degli allenatori (contenuti all'interno della variabile \emph{campionato}) che hanno vinto almeno una coppa.
\end{frame}




\end{document}
